\section{An action estimate that retains the equality defects}
\label{act:section}

The finite-interval argument for the sharp inequality can also measure
failure of its equality relations.  We retain the first-order square
and a quantitative part of the matrix trace deficit.  The coefficient
of this second defect will be independent of the number of rows.  This
uniformity allows us to exhaust an eigenvalue list before knowing that
it is finite.

The construction adapts the field, connecting paths, and action
identity of \cite[Sections~3--6]{MLT}.  Their analytic foundations are
proved in Appendices~\ref{app:trace} and~\ref{app:paths}; the additional
estimate below extracts operator-norm rigidity from the trace deficit.

Fix positive integers $n,d$, a bounded interval
$\Omega=(x_-,x_+)$, and
$K=\operatorname{diag}(\ell_1,\ldots,\ell_n)$ with
$0<\ell_i\leq\kappa$.  For
$A\in H^1_0(\Omega;\R^{n\times d})$, set $\rho=AA^{\mathsf T}$
and define
\begin{equation}\label{act:energy}
 \mathcal E_K(A)=\int_\Omega
 \bigl(\|A'\|_{\HS}^2+\tr(K^2\rho)-\tr(\rho^q)\bigr)\,dx.
\end{equation}
For pointwise matrix values $A\in\R^{n\times d}$ and
$B\in\Sym_n$, put
\begin{equation}\label{act:defect}
 \Phi_K(A,B)=\min\left\{1,
       \|[B,K^2]\|_{\op}^2
       +\|K^2-B^2-(AA^{\mathsf T})^r\|_{\op}^2\right\}.
\end{equation}
Thus $\Phi_K=0$ expresses two algebraic relations, while
$A'-BA=0$ is a first-order differential relation.

\begin{theorem}[Quantitative finite-interval action estimate]
\label{act:quantitative-action}
Let $0<\sigma<1$ and $\kappa>0$.  There is a constant
$c_*=c_*(\sigma,\kappa)>0$ with the following property.
For $n,d,\Omega,K$ as above, suppose
$U\in H^1_0(\Omega;\R^{n\times d})$ satisfies
\[
                  \int_\Omega UU^{\mathsf T}\,dx=I_n.
\]
For every $\eta>0$ there exist a real lower triangular matrix
$C\in\R^{n\times n}$ and a path
$B\in C^1(\overline\Omega;\Sym_n)$ such that
\[
 B(x_-)=K,\qquad B(x_+)=-K,\qquad
 \sup_{x\in\overline\Omega}\|B(x)\|_{\op}\leq\kappa,
\]
and, for $A=CU$,
\begin{equation}\label{act:bound}
 \mathcal E_K(A)\geq b_\sigma\sum_{i=1}^n\ell_i^{2\gamma}
   +\int_\Omega\left(\|A'-BA\|_{\HS}^2
                        +c_*\Phi_K(A,B)\right)\,dx-\eta.
\end{equation}
In particular, $c_*$ is independent of $n,d,\Omega,U$, and the
smallest diagonal entry of $K$.
\end{theorem}

\subsection{Coercivity at the matrix-field constraint}

The algebraic relations $\Phi_K(A,B)=0$ require $[B,K^2]=0$
and $\rho=(K^2-B^2)^\sigma$.  For $0<\delta\leq1$, we use a
locally Lipschitz field $M_\delta:\Sym_n\to\Sym_n$ extending
the commuting rule
\[
 M_\delta(B)=(K^2-B^2+\delta I)^\sigma-\delta^\sigma I
 \quad\text{if }[B,K^2]=0\text{ and }K^2-B^2\geq0.
\]
Thus the constraint $M_\delta(B)=\rho$ is a regularized form of
the algebraic equality relation in this case.  The extension to
arbitrary symmetric $B$ has a $C^1$ primitive
$J_\delta:\Sym_n\to\R$ satisfying
\begin{equation}\label{act:primitive}
 dJ_\delta(B)[H]=-\tr(M_\delta(B)H).
\end{equation}
Proposition~\ref{field:properties} constructs these maps and proves
their required properties.  Along a path $B$, the derivative identity
converts $-\tr(M_\delta(B)B')$ into the total derivative of
$J_\delta(B)$, which will supply the boundary term in the action
identity.
Its endpoint difference is explicit:
\begin{align}
 J_\delta(-K)-J_\delta(K)
 &=\sum_{i=1}^n\int_{-\ell_i}^{\ell_i}
     \bigl((\ell_i^2-s^2+\delta)^\sigma-\delta^\sigma\bigr)\,ds
       \notag\\
 &\longrightarrow b_\sigma\sum_{i=1}^n\ell_i^{2\gamma}
       \qquad(\delta\downarrow0).
       \label{act:endpoints}
\end{align}
The limit follows by dominated convergence and
$2\sigma+1=2\gamma$.

Define the scalar function
\begin{equation}\label{act:F}
 F(A,B)=\tr(\rho^q)-\tr\bigl((K^2-B^2)\rho\bigr),
 \qquad \rho=AA^{\mathsf T}.
\end{equation}
The following lemma bounds this nonlinear term while keeping the
algebraic defect.  The constraint in its hypothesis will later be
enforced by a penalty.

\begin{lemma}\label{act:constraint}
The constant $c_*(\sigma,\kappa)>0$ can be chosen so that, whenever
$0<\delta\leq1$, $\|B\|_{\op}\leq\kappa$, and
$M_\delta(B)=AA^{\mathsf T}$, one has
\begin{equation}\label{act:constraint-bound}
 F(A,B)+c_*\Phi_K(A,B)\leq E_{n,\delta},
 \qquad
 E_{n,\delta}
 =n\bigl(\delta\kappa^{2\sigma}
           +\delta^\sigma(\kappa^2+\delta)\bigr)
      +\delta^{2\sigma}.
\end{equation}
\end{lemma}

\begin{proof}
Write $\rho=M_\delta(B)\geq0$.  Proposition~\ref{field:positive}
gives
\begin{equation}\label{act:positive-symbol}
 Q_s:=s^2I-2sB+K^2\geq0\quad(s\in\R),
 \qquad 0\leq\rho\leq\kappa^{2\sigma}I.
\end{equation}
Put $V=K^2-B^2+\delta I$ and $\beta=\sigma-\tfrac12$.
With the positive normalization constant $c_\sigma$ specified in
\eqref{trace:normalization}, define
\begin{equation}\label{act:RN}
 \begin{split}
 R(t)&=\frac1\pi\int_\R
           \bigl(Q_s+(\delta+t)I\bigr)^{-1}\,ds,\qquad t>0,\\
 N&=c_\sigma\int_0^\infty
           t^\beta\bigl(t^{-1/2}I-R(t)\bigr)\,dt
       =\rho+\delta^\sigma I.
 \end{split}
\end{equation}
The last identity is also part of Proposition~\ref{field:positive}.
Since $(sI-B)^2+V=Q_s+\delta I\geq\delta I$ and $N\geq0$,
Theorem~\ref{trace:deficit} applies.  It supplies
$L(t)=R(t)^{-1}>0$ and Hermitian matrices $Z(t)$ satisfying
\begin{equation}\label{act:LZ}
 V+tI=L^2+Z^2+i[B,Z],\qquad
 LZ+ZL+i[B,L]=0,
\end{equation}
and the exact identity
\begin{align}
 \Delta&:=\tr(NV)-\tr(N^q)\notag\\
 &=c_\sigma\int_0^\infty t^\beta
    \left(\tr(R(t)Z(t)^2)
      +q\bigl\|R(t)^{-1/2}
           -R(t)^{1/2}(tI+N^r)^{1/2}\bigr\|_{\HS}^2\right)\,dt
       \geq0.
       \label{act:trace-deficit}
\end{align}
The matrices in \eqref{act:LZ} may be complex Hermitian even though
$A$ and $B$ are real; $i$ denotes the imaginary unit.

We extract a dimension-independent bound from the part of the
integral with $1\leq t\leq2$.  On this interval there are constants
$a_\kappa,A_\kappa>0$ such that
\begin{equation}\label{act:R-bounds}
                       a_\kappa I\leq R(t)\leq A_\kappa I.
\end{equation}
For the upper bound, use $Q_s\geq0$ on a bounded interval in $s$,
and $Q_s\geq\tfrac12s^2I$ for $|s|\geq\max\{1,4\kappa\}$.
For the lower bound, restrict the integral to $|s|\leq1$, where
$Q_s+(\delta+t)I\leq(4+2\kappa+\kappa^2)I$.
These estimates depend on $\kappa$ alone.  Moreover
$\|N\|_{\op}\leq\kappa^{2\sigma}+1$, so
$H_t=(tI+N^r)^{1/2}$ and $L(t)$ are bounded in operator norm by
constants depending only on $\sigma,\kappa$.

The two integrands in \eqref{act:trace-deficit} control
$\|Z(t)\|_{\HS}^2$ and $\|L(t)-H_t\|_{\HS}^2$:
indeed $\tr(RZ^2)=\tr(ZRZ)$ and
\[
 R^{-1/2}-R^{1/2}H_t=R^{1/2}(L-H_t).
\]
It follows by averaging over $[1,2]$ that at some $t$ in this
interval,
\begin{equation}\label{act:small-squares}
 \|Z(t)\|_{\op}+\|L(t)-H_t\|_{\op}
        \leq C_1\sqrt\Delta,
\end{equation}
where $C_1$ depends only on $\sigma,\kappa$.
When $\Delta=0$, the same conclusion follows by taking a point
where both nonnegative integrands vanish.

Suppose first that $\Delta\leq1$, and use this value of $t$.
The first identity in \eqref{act:LZ} gives
\[
 V-N^r=(L^2-H_t^2)+Z^2+i[B,Z].
\]
The boundedness of $L,H_t,B$, together with
\eqref{act:small-squares}, bounds the norm of this expression by
$C_2\sqrt\Delta$.  To control the commutator, the second identity
in \eqref{act:LZ} first yields
$\|[B,L]\|_{\op}\leq2\|L\|_{\op}\|Z\|_{\op}$.
Then
\[
 [B,V]=[B,L^2]+[B,Z^2]+i[B,[B,Z]]
\]
has the same bound.  Thus
\begin{equation}\label{act:V-coercivity}
 \|V-N^r\|_{\op}+\|[B,V]\|_{\op}
       \leq C_3\sqrt\Delta.
\end{equation}
Here and below the constants in this proof depend only on
$\sigma,\kappa$.

The matrices $N=\rho+\delta^\sigma I$ and $\rho$ commute.
Because $r>1$, the scalar function $y\mapsto y^r$ is Lipschitz on
$[0,\kappa^{2\sigma}+1]$.  Consequently
\[
 \|N^r-\rho^r\|_{\op}\leq C_4\delta^\sigma.
\]
Since $[B,V]=[B,K^2]$ and $\delta\leq\delta^\sigma$,
\eqref{act:V-coercivity} implies
\[
 \|[B,K^2]\|_{\op}^2
   +\|K^2-B^2-\rho^r\|_{\op}^2
       \leq C_5(\Delta+\delta^{2\sigma}).
\]
Choose $c_*>0$ sufficiently small, also with $c_*\leq1$.
For $\Delta\leq1$ the preceding estimate gives
\begin{equation}\label{act:coercivity}
                   c_*\Phi_K(A,B)\leq\Delta+\delta^{2\sigma}.
\end{equation}
For $\Delta>1$ this follows directly from $\Phi_K\leq1$.
This proves the required uniform control of the algebraic defect.

It remains to relate $\Delta$ to $F$.  Expanding $N$ and $V$ gives
\[
 F(A,B)=-\Delta+\tr(\rho^q)-\tr(N^q)
                  +\delta\tr\rho+\delta^\sigma\tr V.
\]
The first two powers are ordered because $N=\rho+\delta^\sigma I$.
Also $\tr\rho\leq n\kappa^{2\sigma}$ and
$V\leq(\kappa^2+\delta)I$.  Adding \eqref{act:coercivity} now
proves \eqref{act:constraint-bound}.
\end{proof}

\subsection{The action identity and the penalty limit}

\begin{proof}[Proof of Theorem~\ref{act:quantitative-action}]
Fix $0<\delta\leq1$.  Proposition~\ref{path:connecting} supplies,
for every $0<\varepsilon\leq1$, a lower triangular matrix
$C_\varepsilon$ and a path $B_\varepsilon$ with the prescribed
endpoints and operator-norm bound such that, for
$A_\varepsilon=C_\varepsilon U$,
\begin{equation}\label{act:path}
 \varepsilon B_\varepsilon'
       =M_\delta(B_\varepsilon)
                    -A_\varepsilon A_\varepsilon^{\mathsf T}.
\end{equation}
For fixed $\delta,K,\Omega,U$, the matrices $C_\varepsilon$ are
bounded uniformly in $0<\varepsilon\leq1$.

Temporarily suppress the subscript $\varepsilon$ and write
\[
                 u_\delta(A,B)
                      =\|M_\delta(B)-AA^{\mathsf T}\|_{\HS}^2.
\]
Since $A\in H^1_0$ and $B\in C^1$, the product and chain rules
for absolutely continuous functions, \eqref{act:primitive}, and
\eqref{act:path} yield almost everywhere
\begin{align*}
 \bigl(J_\delta(B)+\tr(A^{\mathsf T}BA)\bigr)'
 &=2\tr((A')^{\mathsf T}BA)
          -\tr\bigl((M_\delta(B)-AA^{\mathsf T})B'\bigr)\\
 &=2\tr((A')^{\mathsf T}BA)-\varepsilon^{-1}u_\delta(A,B).
\end{align*}
Complete the square in the action and integrate this identity.
Because $A$ vanishes at both endpoints, the result is the exact
formula
\begin{align}
 \mathcal E_K(A)
 &=J_\delta(-K)-J_\delta(K)\notag\\
 &\quad+\int_\Omega\left(
       \|A'-BA\|_{\HS}^2-F(A,B)
                          +\varepsilon^{-1}u_\delta(A,B)\right)\,dx.
       \label{act:identity}
\end{align}

We now use the penalty to pass from the constraint estimate in
Lemma~\ref{act:constraint} to all matrix values on these paths.
The continuous representative of $U$ is bounded on
$\overline\Omega$.  The uniform bound on $C_\varepsilon$ and the
bound $\|B_\varepsilon\|_{\op}\leq\kappa$ therefore place all pairs
$(A_\varepsilon(x),B_\varepsilon(x))$ in one compact set
$\mathcal K_\delta\subset\R^{n\times d}\times\Sym_n$ on which
$\|B\|_{\op}\leq\kappa$.  On this set $F$, $\Phi_K$, and
$u_\delta$ are continuous.  Lemma~\ref{act:constraint} implies
\begin{equation}\label{act:penalty}
 \limsup_{\varepsilon\downarrow0}
 \sup_{(A,B)\in\mathcal K_\delta}
    \bigl(F(A,B)+c_*\Phi_K(A,B)-\varepsilon^{-1}u_\delta(A,B)\bigr)
       \leq E_{n,\delta}.
\end{equation}
Indeed, if the left side exceeded $E_{n,\delta}$, some sequence
$\varepsilon_j\downarrow0$ and corresponding pairs in the compact
set would make the expression at least $E_{n,\delta}+\tau$ for
some $\tau>0$.  Boundedness of $F+c_*\Phi_K$ would force
$u_\delta=O(\varepsilon_j)$ along that sequence.  A convergent
subsequence would then give a limiting pair with $u_\delta=0$ and
$F+c_*\Phi_K\geq E_{n,\delta}+\tau$, contradicting
Lemma~\ref{act:constraint}.

For any $\tau>0$, choose $\varepsilon$ small enough that the
supremum in \eqref{act:penalty} is at most $E_{n,\delta}+\tau$.
Substitution into \eqref{act:identity} gives
\begin{align*}
 \mathcal E_K(A_\varepsilon)
 \geq{}&J_\delta(-K)-J_\delta(K)
        -|\Omega|(E_{n,\delta}+\tau)\\
 &+\int_\Omega\left(\|A_\varepsilon'-B_\varepsilon A_\varepsilon\|_{\HS}^2
                +c_*\Phi_K(A_\varepsilon,B_\varepsilon)\right)\,dx.
\end{align*}
Finally, first choose $\delta>0$ so small that the endpoint error
in \eqref{act:endpoints} and $|\Omega|E_{n,\delta}$ together are
less than $\eta/2$.  Next choose $\tau>0$ with
$|\Omega|\tau<\eta/2$, and then choose the required $\varepsilon$.
These choices prove \eqref{act:bound}.  All compactness arguments
were made after fixing $\delta$; no bound uniform in $\delta$ is
needed.
\end{proof}
